\section{Auto-encoding variational Bayes}

Auto-encoding variational Bayes method \cite{kingma2022autoencodingvariationalbayes} is designed for problems with latent variables and intractable distributions so that 
EM and variational inference do not work. 
The method uses generative models to create variational posterior and optimize model parameters by maximizing Evidence Lower Bound (ELBO). 
Hence it provides a framework for combining neural networks and Bayesian inference. 

\subsection{Preliminary}
\subsubsection{Monte Carlo gradient estimator }

Suppose we have a random variable $z\sim q_{\phi}(z)$ where $\phi$ is the parameter of the density $q$. 
For an arbitrary function $f\in L_2(q_{\phi})$, 
we want to estimate the gradient of $\mathbb{E}_{q_{\phi}(z)}[f(z)]$ w.r.t $\phi$. 
We have two approaches that use Monte Carlo method to approximate the expectation. 

\begin{equation}
    \label{eq: sec1_MCGE1}
    \nabla_{\phi}\mathbb{E}_{q_{\phi}(z)}[f(z)] = \mathbb{E}_{q_{\phi}(z)}[f(z)\nabla_{\phi}\log q_{\phi}(z)]\approx \frac{1}{N}\sum_{i=1}^{N}f(z_i)\nabla_{\phi}\log q_{\phi}(z_i).
\end{equation}

Suppose $z = g(\phi,\epsilon)$, $\epsilon\sim p(\epsilon)$. Then $p(\epsilon)=q_{\phi}(g(\phi,\epsilon))\frac{\partial g(\phi,\epsilon)}{\partial \epsilon}$. 
\begin{equation}
    \label{eq: sec1_MCGE2}
    \nabla_{\phi}\mathbb{E}_{q_{\phi}(z)}[f(z)] = \mathbb{E}_{\epsilon}[\nabla_{\phi}f(g(\phi,\epsilon))] \approx \frac{1}{N}\sum_{i=1}^{N}\nabla_{\phi}f(g(\phi,\epsilon_i)).
\end{equation}

Method (\ref{eq: sec1_MCGE1}) uses score function $\nabla_{\phi}\log q_{\phi}(z)$ and (\ref{eq: sec1_MCGE2}) uses reparameterization. 
The score-function approach is unbiased but with high variance. 
The reparameterization approach has low variance but requires a differentiable pathway. 

\subsubsection{Evidence lower bound (ELBO)}

Suppose the observable variable $X$ and the latent variable $Z$ belong to the same parametric family $p_{\theta}$, $\theta\in\Theta\subseteq\mathbb{R}^d$. 
To estimate the true posterior $p_{\theta}(z|x)$, we use maximum likelihood method to optimize the parameter $\theta$.
However when the true marginal density $p_{\theta}(x)=\int p_{\theta}(x|z)p_{\theta}(z)dz$ is intractable, 
by Bayes formula, $p_{\theta}(z|x)=\frac{p_{\theta}(x|z)p_{\theta}(z)}{p_{\theta}(x)}$ is also intractable. 
This makes the gradient optimizer lose its power. 
One optional approach is to use variational density $q_{\phi}(z|x)$ to approximate the true posterior. 
And the variational posterior is optimized by maximizing the ELBO.  

\begin{align}
    \label{eq: sec1_ELBO1}
    \log p_{\theta}(x) &= \int\log p_{\theta}(x)q_{\phi}(z|x)dz \notag\\
                       &= \int(\log p_{\theta}(z,x) - \log p_{\theta}(z|x))q_{\phi}(z|x)dz \notag\\
                       &= \int(\log p_{\theta}(z,x) - \log q_{\phi}(z|x))q_{\phi}(z|x)dz + \int(\log q_{\phi}(z|x) - \log p_{\theta}(z|x))q_{\phi}(z|x)dz \notag\\
                       &= \mathcal{L}(\theta,\phi;x) + D_{KL}(q_{\phi}(z|x)||p_{\theta}(z|x)).
\end{align}

By (\ref{eq: sec1_ELBO1}), since Kullback-Leibler divergence is nonnegative and $\log p_{\theta}(x)$ is a constant irrelevant to $\phi$, 
we can maximize the ELBO $\mathcal{L}(\theta,\phi;x)$ to minimize the KL divergence between the variational posterior and the true posterior. 

Another expression of ELBO is given by (\ref{eq: sec1_ELBO2}).

\begin{align}
    \label{eq: sec1_ELBO2}
    \mathcal{L}(\theta,\phi;x) &= \int(\log p_{\theta}(z,x) - \log q_{\phi}(z|x))q_{\phi}(z|x)dz \notag\\
                               &= \int(\log p_{\theta}(x|z) + \log p_{\theta}(z) - \log q_{\phi}(z|x))q_{\phi}(z|x)dz \notag\\
                               &= \int\log p_{\theta}(x|z)q_{\phi}(z|x)dz - \int(\log q_{\phi}(z|x) - \log p_{\theta}(z))q_{\phi}(z|x)dz \notag\\
                               &= \mathbb{E}_{q_{\phi}(z|x)}[\log p_{\theta}(x|z)] - D_{KL}(q_{\phi}(z|x)||p_{\theta}(z)).
\end{align}

\subsection{Method}

The Auto-encoding variational Bayes algorithm uses arbitrary encoder to encode data to posterior function (parameters).  
The latent state (embedding) is sampled from the variational posterior. 
The model learnable parameters are updated using the gradient computed from the ELBO. 
The algorithm not only gives freedom to choose encoder and decoder but also provides opportunity to combine other deterministic methods to embed data to latent state variable. 

\subsubsection{Auto-encoding variational Bayes algorithm}

From (\ref{eq: sec1_ELBO1}) and (\ref{eq: sec1_ELBO2}), we have two expressions for the ELBO. 
\begin{equation}
    \label{eq: sec1_ELBO3}
    \mathcal{L}(\theta,\phi;x) = \left\{
    \begin{aligned}
        & \mathbb{E}_{q_{\phi}(z|x)}[\log p_{\theta}(z,x) - \log q_{\phi}(z|x)] \\
        & \mathbb{E}_{q_{\phi}(z|x)}[\log p_{\theta}(x|z)] - D_{KL}(q_{\phi}(z|x)||p_{\theta}(z)) 
    \end{aligned}\right.
\end{equation}
For any differentiable pathway $g_{\phi}$ and auxiliary noise variable $\epsilon\sim p(\epsilon)$, 
we apply the reparameterization trick to (\ref{eq: sec1_ELBO3}) to get the SGVBs. 

For an observation $X=x$, we let $Z = g_{\phi}(\epsilon, x)$. 
Then the i.i.d. Monte Carlo sample $z_1,\cdots,z_N\sim q_{\phi}(z|x)$ are determined by i.i.d. $\epsilon_1,\cdots,\epsilon_N\sim p(\epsilon)$. 
And the ELBO is approximated by 
\begin{equation}
    \tilde{\mathcal{L}}(\theta,\phi;x) = \left\{
    \begin{aligned}
        & \frac{1}{N}\sum_{i=1}^{N}(\log p_{\theta}(z_i,x) - \log q_{\phi}(z_i|x)) \\
        & \frac{1}{N}\sum_{i=1}^{N}\log p_{\theta}(x|z_i) - D_{KL}(q_{\phi}(z_i|x)||p_{\theta}(z_i)) 
    \end{aligned}\right.
\end{equation}
One can take any of the formations to approximate the ELBO. 
Since the KL divergence usually can be analytically computed, we don't need to approximate it with Monte Carlo method. 

The Auto-encoding variational Bayes algorithm (Minibatch version) is given as follows. 

% \hrulefill

% \textbf{The Auto-encoding variational Bayes algorithm}

% \hrulefill

% Initialize learnable parameters $\theta^{(0)}$ 

% Define encoder $G$

% Impose latent variable prior $p(Z)$

% \textbf{For k=1:K}

% \quad i) $\phi^{(k)}\longleftarrow$ Encode data to posterior parameter: $\phi^{(k)}=G_{\theta^{(k)}}(X)$

% \quad ii) $z_{k}\longleftarrow$ Sample latent state (embedding) from posterior using reparameterization $g_{\phi^{(k)}}$

%     \qquad iia) $\epsilon\longleftarrow$ Random samples from noise distribution $p(\epsilon)$

%     \qquad iib) $z_{k}\longleftarrow g_{\phi^{(k)}}(\epsilon)$ Derive the sampled latent variable

% \quad iii) $\nabla_{\theta^{(k)}}\tilde{\mathcal{L}}(\theta^{(k)};z_{k},X)$ Compute gradients w.r.t model parameters 

% \quad iv) $\theta^{(k+1)}\longleftarrow$ Update model parameters using gradients

% \quad \textbf{If Converge then break}

% \textbf{Return final parameters}

% \hrulefill

\begin{algorithm}
    \caption{Auto-encoding Variational Bayes}
    \begin{algorithmic}
        \State Initialize $\theta^{(0)}$
        \For{$k=1:K$}
        \State $\phi^{(k)} \gets G_{\theta^{(k)}}(X)$
        \State sample $\epsilon \sim p(\epsilon)$
        \State $z_k \gets g_{\phi^{(k)}}(\epsilon)$
        \State compute gradient
        \State update $\theta$
        \EndFor
        \State \Return parameters
    \end{algorithmic}
\end{algorithm}

\subsubsection{Variational Auto-encoders}

The Auto-encoding variational Bayes algorithm provides freedom for choosing the variational posterior $q_{\phi}(z|x)$. 
When the variational posterior is determined by a neural network, 
the algorithm is equivalent to an auto-encoder which uses a probabilistic encoder $q_{\phi}(z|x)$. 

As an example, we assume the posterior belongs to the Gaussian family, i.e. $q_{\phi}(z|x)=\mathcal{N}(\mu_x,\sigma_x^2)$. 
The parameters $\mu_x$ and $\sigma_x^2$ are mapped from observation $x$ by MLPs with parameters $\phi=(W_0,W_1,W_2,b_0,b_1,b_2)$. 
This process is defined as follows. 
\begin{align*}
    h_x &= tanh(W_0x+b_0) \\
    \mu_x &= W_1h_x+b_1 \\
    \log\sigma_x^2 &= W_2h_x+b_2
\end{align*}
Let the auxiliary random variable $\epsilon\sim p(\epsilon)=\mathcal{N}(0,1)$. 
Then the differentiable pathway $g_{\phi}$ is given by 
\begin{equation*}
    \forall z\sim q_{\phi}(z|x),\ \exists\epsilon\sim p(\epsilon),\ s.t.\ z=g_{\phi}(x,\epsilon)=\mu_x + \sigma_x^2\epsilon. 
\end{equation*}
The probabilistic decoding process is also determined by a MLPs with the same structure but different hyperparameters $\theta=(w_3,w_4,w_5,b_3,b_4,b_5)$. 
\begin{align*}
    h_z &= tanh(W_3z+b_3) \\
    \mu_z &= W_4h_z+b_4 \\
    \log\sigma_z^2 &= W_5h_z+b_5
\end{align*}
We sample the prediction $x_{pred}$ from $p_{\theta}(x|z)=\mathcal{N}(\mu_z,\sigma_z^2)$. 

The non-probabilistic decoder can be an activation function operating on the inner product of embeddings. 
\begin{equation*}
    x_{pred} = \sigma(\left\langle z, z\right\rangle).
\end{equation*}